\documentclass[10pt,a4paper]{amsart}
\usepackage[margin=19mm]{geometry}
\usepackage{amsmath,amssymb,mathtools}
\usepackage[scr=rsfs,cal=euler]{mathalfa}
\usepackage{xfrac}
\usepackage[unicode,hidelinks,bookmarks=false]{hyperref}
\newcommand{\ie}{i.e.\ }
\newcommand{\cf}{cf.\ }
\newcommand{\resp}{respectively\ }
\newcommand{\Subsection}{\subsection}
\providecommand{\nobreakdash}{\nobreak\hbox{-}}
\setlength{\emergencystretch}{2em}
\multlinegap0pt
\usepackage{amssymb}
\usepackage{xcolor}
\usepackage{float}
\newtheorem{lemma}{Lemma}
\newtheorem{proposition}{Proposition}
\usepackage{cleveref}
\crefname{lemma}{Lemma}{Lemmas}
\crefname{proposition}{Proposition}{Propositions}
\newcommand{\E}{\mathbb{E}}
\title{The Random Subsequence Model and Uniform Codes for the Deletion Channel}
\author{Ryan Jeong, Francisco Pernice}
\date{Source: arXiv:2604.07234v1 (CC BY 4.0)}
\begin{document}
\maketitle

\noindent\emph{Source and licence.} The Random Subsequence Model and Uniform Codes for the Deletion Channel. Version arXiv:2604.07234v1 (CC BY 4.0), distributed by arXiv under the Creative Commons Attribution 4.0 International licence, http://creativecommons.org/licenses/by/4.0/, as recorded in the arXiv licence metadata for this exact version and shown on the abstract page https://arxiv.org/abs/2604.07234v1. Full licence notice: \texttt{licenses/arxiv-2604.07234-CC-BY-4.0.txt}.\par\smallskip

\noindent\textit{Editorial scope.} Counted unit: the article's proposition prop:Z-equals-one (source0.tex lines 1880--1895). The article's complete original proof of it is delivered with it (source0.tex lines 1897--2042, one \texttt{proof} environment of 146 lines). No mathematics is written, repaired or re-derived: the statement and the proof are the article's own lines, in the article's own notation, and no retained line is substituted or re-flowed.  Retained uncounted as the source-local context the proof needs: lines 1735--1748; lines 1745--1747; lines 1814--1820.  Nothing was omitted from any retained span, so the package carries no omission record.  No retained line contains a citation, so the wrapper carries no bibliography and no bibliography entry.  No retained line contains a figure environment, an image-inclusion command or drawing code, so no image is delivered and the package declares no asset dependency.  Editorial formatting only, in addition: an AMS \texttt{amsart} preamble that restates the article's own declarations and macros used by the retained lines, namely the environments \texttt{lemma}, \texttt{proposition} on separate counters, together with the standard packages those lines need.  The retained environments print in this document as Lemma 1, Proposition 1 in order of appearance, whereas the article prints the same items with the numbering of its own full document.  The complete native source of the article as distributed in the arXiv e-print for this exact version is bundled unchanged as \texttt{sources/198077/source0.tex}.
\begin{lemma}[Exponential family optimizer]\label{lem:inner-optimizer}
Fix $\rho\in(0,1).$ The supremum defining $\Phi(\rho)$ is attained by a unique probability measure $\nu^*$ of the form
\begin{equation}\label{eq:inner-optimizer-form}
    \nu^*(a,b)\;=\;\frac{w(a,b)\,x^a\,y^b}{Z(x,y)},\qquad a\ge 1,\;1\le b\le a,
\end{equation}
where $x,y>0$ are the unique positive solutions of the moment equations $\E_{\nu^*}[a]=1/(\alpha\rho)$ and $\E_{\nu^*}[b]=1/\rho$, and
\begin{align*}
    Z(x,y)\;:=\;\sum_{a\ge 1}\sum_{b=1}^{a}\binom{a-1}{b-1}^{2}\,x^{a}\,y^{b}
\end{align*}
is the partition function. The optimal value is
\begin{equation}\label{eq:Phi-from-Z}
    \Phi(\rho)\;=\;\log Z(x,y)\;-\;\frac{1}{\alpha\rho}\,\log x\;-\;\frac{1}{\rho}\,\log y.
\end{equation}
\end{lemma}

\begin{equation}
    \Phi(\rho)\;=\;\log Z(x,y)\;-\;\frac{1}{\alpha\rho}\,\log x\;-\;\frac{1}{\rho}\,\log y.
\end{equation}

\begin{lemma}[Closed-form partition function]\label{lem:partition-closed-form}
    For all $x,y>0$ with $(1-x-xy)^2>4x^2 y$,
    \begin{equation}\label{eq:partition-closed-form}
        Z(x,y)\;=\;\frac{xy}{\sqrt{(1-x-xy)^{2}-4x^{2}y}}.
    \end{equation}
    Moreover, $Z(x,y) = \infty$ for all $x, y > 0$ that fail to satisfy this condition.
\end{lemma}

\begin{proposition}[Unique interior maximizer and normalization]\label{prop:Z-equals-one}
Let
\[
g(\rho):=\alpha\rho\,\Phi(\rho), \qquad \rho\in(0,1).
\]
Then \(g\) attains its maximum at a unique point \(\rho^*\in(0,1)\). If
\((x^*,y^*)\) denotes the pair of parameters from \cref{lem:inner-optimizer}
corresponding to \(\rho=\rho^*\), then
\begin{equation}\label{eq:Z-equals-one}
    Z(x^*,y^*)=1,
\end{equation}
and consequently
\begin{equation}\label{eq:R-explicit}
    \mathsf R(\alpha)\;=\;-\log x^*\;-\;\alpha\log y^*.
\end{equation}
\end{proposition}

\begin{proof}
For each \(\rho\in(0,1)\), let \(x=x(\rho)\) and \(y=y(\rho)\) be the unique
positive parameters given by \cref{lem:inner-optimizer}. By \eqref{eq:Phi-from-Z},
\begin{equation}\label{eq:g-rho-from-Z}
    g(\rho)=\alpha\rho\log Z(x,y)-\log x-\alpha\log y.
\end{equation}

We first compute \(g'(\rho)\). The Lagrange multipliers for the inner problem are
\[
\alpha_1=-\log x
\qquad\text{and}\qquad
\alpha_2=-\log y,
\]
corresponding respectively to the constraints
\[
\E_{\nu^*}[a]=\frac{1}{\alpha\rho},
\qquad
\E_{\nu^*}[b]=\frac{1}{\rho}.
\]
By the envelope theorem,
\[
\frac{d\Phi}{d\rho}
=
\alpha_1\frac{d(1/(\alpha\rho))}{d\rho}
+
\alpha_2\frac{d(1/\rho)}{d\rho}
=
\frac{\frac1\alpha\log x+\log y}{\rho^2}.
\]
Therefore,
\begin{align}
g'(\rho)
&=
\alpha\Phi(\rho)+\alpha\rho\,\Phi'(\rho) \notag\\
&=
\alpha\Bigl[\log Z-\frac{\frac1\alpha\log x+\log y}{\rho}\Bigr]
+
\alpha\frac{\frac1\alpha\log x+\log y}{\rho}
=
\alpha\log Z(x(\rho),y(\rho)).
\label{eq:g-prime-alpha-logZ}
\end{align}

Next we compute \(x(\rho)\), \(y(\rho)\), and \(Z(x(\rho),y(\rho))\) explicitly.
By \cref{lem:partition-closed-form},
\[
Z(x,y)=\frac{xy}{\sqrt{D(x,y)}},
\qquad
D(x,y):=(1-x-xy)^2-4x^2y.
\]
Hence
\[
\log Z(x,y)=\log x+\log y-\frac12\log D(x,y).
\]
On the other hand, by the exponential family form of $\nu^*,$ we have
\[
\E_{\nu^*}[a]=x\partial_x\log Z(x,y),
\qquad
\E_{\nu^*}[b]=y\partial_y\log Z(x,y).
\]
A direct computation gives
\[
x\partial_x\log Z(x,y)=\frac{1-x-xy}{D(x,y)},
\qquad
y\partial_y\log Z(x,y)=\frac{1-2x-xy+x^2(1-y)}{D(x,y)}.
\]
Therefore the moment constraints are equivalent to
\begin{equation}\label{eq:moment-system-a}
    \frac{1-x-xy}{D(x,y)}=\frac{1}{\alpha\rho},
\end{equation}
and
\begin{equation}\label{eq:moment-system-b}
    \frac{1-2x-xy+x^2(1-y)}{D(x,y)}=\frac{1}{\rho}.
\end{equation}
Solving \eqref{eq:moment-system-a}--\eqref{eq:moment-system-b} yields
\begin{equation}\label{eq:x-rho-explicit}
    x(\rho)=\frac{(1-\alpha)(2-2\alpha+\alpha\rho)}{2-\alpha\rho},
\end{equation}
and
\begin{equation}\label{eq:y-rho-explicit}
    y(\rho)=\frac{\alpha^2(2-\rho)(1-\rho)}{(1-\alpha)(2-2\alpha+\alpha\rho)}.
\end{equation}
Substituting \eqref{eq:x-rho-explicit}--\eqref{eq:y-rho-explicit} into the
closed form for \(Z\) gives
\begin{equation}\label{eq:Z-rho-explicit}
    Z(x(\rho),y(\rho))
    =
    \frac{\alpha(1-\rho)\sqrt{2-\rho}}
    {\sqrt{\rho}\,\sqrt{(2-\alpha\rho)(2-2\alpha+\alpha\rho)}}.
\end{equation}

In particular,
\[
Z(x(\rho),y(\rho))\longrightarrow\infty
\qquad\text{as }\rho\downarrow 0,
\]
and
\[
Z(x(\rho),y(\rho))\longrightarrow 0
\qquad\text{as }\rho\uparrow 1.
\]
Moreover, \(Z(x(\rho),y(\rho))=1\) is equivalent, after squaring
\eqref{eq:Z-rho-explicit}, to
\begin{equation}\label{eq:rho-quadratic}
    2\alpha^2\rho^2+(4\alpha-5\alpha^2-4)\rho+2\alpha^2=0.
\end{equation}
The polynomial on the left-hand side of \eqref{eq:rho-quadratic} has value
\(2\alpha^2>0\) at \(\rho=0\) and value
\[
2\alpha^2+(4\alpha-5\alpha^2-4)+2\alpha^2
=
4\alpha-\alpha^2-4
=
-(2-\alpha)^2<0
\]
at \(\rho=1\). Hence it has at least one root in \((0,1)\). Since its constant
term equals its leading coefficient, the product of its roots is \(1\), so it
can have at most one root in \((0,1)\). Therefore there is a unique
\(\rho^*\in(0,1)\) such that
\[
Z(x(\rho^*),y(\rho^*))=1.
\]

Because \(Z(x(\rho),y(\rho))\) is continuous, the uniqueness of \(\rho^*\) and
the above boundary limits imply
\[
Z(x(\rho),y(\rho))>1 \quad \text{for } 0<\rho<\rho^*,
\qquad
Z(x(\rho),y(\rho))<1 \quad \text{for } \rho^*<\rho<1.
\]
By \eqref{eq:g-prime-alpha-logZ}, it follows that
\[
g'(\rho)>0 \quad \text{for } 0<\rho<\rho^*,
\qquad
g'(\rho)<0 \quad \text{for } \rho^*<\rho<1.
\]
Thus \(g\) is strictly increasing on \((0,\rho^*)\) and strictly decreasing on
\((\rho^*,1)\), so \(g\) attains its maximum uniquely at \(\rho^*\in(0,1)\).

Finally, setting \((x^*,y^*):=(x(\rho^*),y(\rho^*))\), we have proved
\eqref{eq:Z-equals-one}. Plugging this into \eqref{eq:g-rho-from-Z} gives
\[
\mathsf R(\alpha)=g(\rho^*)=-\log x^*-\alpha\log y^*,
\]
which is \eqref{eq:R-explicit}.
\end{proof}

\end{document}
